\documentclass[11pt,a4paper]{article}
\usepackage[T1]{fontenc}
\usepackage[utf8]{inputenc}
\usepackage{amsmath,amssymb,amsthm}
\usepackage[margin=1in]{geometry}
\usepackage[colorlinks=true,linkcolor=blue,urlcolor=blue]{hyperref}
\theoremstyle{plain}
\newtheorem{theorem}{Theorem}
\newtheorem{lemma}[theorem]{Lemma}
\newtheorem{proposition}[theorem]{Proposition}
\newtheorem{corollary}[theorem]{Corollary}
\theoremstyle{definition}
\newtheorem{remark}[theorem]{Remark}

\title{The empirical closed form for OEIS A227331, proved}
\author{Adrian Perez Fontelles\\ \small Independent researcher}
\date{1 September 2026}
\begin{document}
\maketitle

\begin{abstract}
OEIS A227331 carries an empirical closed form for $a(n)$ --- not a recurrence relating
consecutive terms but an explicit formula. It is true, and it is decidable rather than
empirical. The entry counts arrays subject to a condition on a bounded window of consecutive lines, so the lines of an array are the
vertices of a finite digraph, an array is a walk in it, and $a(n)$ is a walk count on
$S=2745$ vertices. Any function of the form $\sum_j p_j(n)\lambda_j^{\,n}$ satisfies the
monic linear recurrence whose characteristic polynomial is
$\prod_j(x-\lambda_j)^{\deg p_j+1}$; here that polynomial is $q(x)=(x-1)^{32}$, of degree
$32$. Two things are then checked exactly: that the walk count satisfies the same
recurrence from some index on, which the Cayley--Hamilton theorem reduces to a finite
computation, and that the two sequences agree at $32$ consecutive indices past that
point. A monic recurrence determines every later term from its predecessors, so agreement
there is agreement for good.
\end{abstract}

\noindent\small 2020 Mathematics Subject Classification. 05A15, 11B37, 15A18.\normalsize

\section{The conjecture}

OEIS A227331 is ``Number of n X 4 0,1 arrays indicating 2 X 2 subblocks of some larger (n+1) X 5 binary array having two adjacent 1's and two adjacent 0's, with rows and columns of the latter in lexicographically nondecreasing order.''. It has offset $1$ and begins
\[
15, 104, 708, 4917, 33297, 204206, 1108382, 5361624,\ \dots
\]
The entry states, as a conjecture contributed by its author and never marked settled:
\begin{quote}
Empirical polynomial of degree 31 (see link above).

Empirical: a(n) = (1/8222838654177922817725562880000000)*n\^{}31 + (1/265252859812191058636308480000000)*n\^{}30 + (211/53050571962438211727261696000000)*n\^{}29 - (43/1829330067670283163009024000000)*n\^{}28 + (22651/435554778016734086430720000000)*n\^{}27 - (266333/145184926005578028810240000000)*n\^{}26 + (211/524821249466102395699200)*n\^{}25 - (10684127/493746307063504227532800000)*n\^{}24 + (101609721163/46905899171032901615616000000)*n\^{}23 - (256487176249/2039386920479691374592000000)*n\^{}22 + (53972039063/6474244191999020236800000)*n\^{}21 - (8512773014233/19422732575997060710400000)*n\^{}20 + (587506424475709061/26512029966235987869696000000)*n\^{}19 - (1364203300425262771/1395369998222946729984000000)*n\^{}18 + (10991879100452675147/279073999644589345996800000)*n\^{}17 - (23226625879260075527/16416117626152314470400000)*n\^{}16 + (8616722338495316891/188767586339684352000000)*n\^{}15 - (21806075561997009195587/16611547597892222976000000)*n\^{}14 + (1283922986480361672788173/38238504758994213273600000)*n\^{}13 - (29049282179539555567087271/38238504758994213273600000)*n\^{}12 + (7228081086426837457050225271/477981309487427665920000000)*n\^{}11 - (62854978468803534961683372187/238990654743713832960000000)*n\^{}10 + (51877457786090129705786443241/13087583474060519424000000)*n\^{}9 - (149063481533132633801271084421/2908351883124559872000000)*n\^{}8 + (61921208895254661989988212567311/110275008901806228480000000)*n\^{}7 - (9925722646170490600108095270877/1934649278979056640000000)*n\^{}6 + (47393917262283251844887075621/1237654420895692800000)*n\^{}5 - (1390548650132864765781401508097/6126389383433679360000)*n\^{}4 + (433846309064646404522540308733/423012600284706432000)*n\^{}3 - (926612557000804897407362797/279770238283536000)*n\^{}2 + (49062604628325366151963/7220177644680)*n - 6644041492 for n>10
\end{quote}
As of the ``Last modified'' line on the live entry (Oct 03 2025 23:58:04, revision 9) this is still
recorded as empirical, and nothing on the entry records it as proved. Write
\[
f(n)\;=\;\frac{n^{31}}{8222838654177922817725562880000000} + \frac{n^{30}}{265252859812191058636308480000000} + \frac{211 n^{29}}{53050571962438211727261696000000} - \frac{43 n^{28}}{1829330067670283163009024000000} + \frac{22651 n^{27}}{435554778016734086430720000000} - \frac{266333 n^{26}}{145184926005578028810240000000} + \frac{211 n^{25}}{524821249466102395699200} - \frac{10684127 n^{24}}{493746307063504227532800000} + \frac{101609721163 n^{23}}{46905899171032901615616000000} - \frac{256487176249 n^{22}}{2039386920479691374592000000} + \frac{53972039063 n^{21}}{6474244191999020236800000} - \frac{8512773014233 n^{20}}{19422732575997060710400000} + \frac{587506424475709061 n^{19}}{26512029966235987869696000000} - \frac{1364203300425262771 n^{18}}{1395369998222946729984000000} + \frac{10991879100452675147 n^{17}}{279073999644589345996800000} - \frac{23226625879260075527 n^{16}}{16416117626152314470400000} + \frac{8616722338495316891 n^{15}}{188767586339684352000000} - \frac{21806075561997009195587 n^{14}}{16611547597892222976000000} + \frac{1283922986480361672788173 n^{13}}{38238504758994213273600000} - \frac{29049282179539555567087271 n^{12}}{38238504758994213273600000} + \frac{7228081086426837457050225271 n^{11}}{477981309487427665920000000} - \frac{62854978468803534961683372187 n^{10}}{238990654743713832960000000} + \frac{51877457786090129705786443241 n^{9}}{13087583474060519424000000} - \frac{149063481533132633801271084421 n^{8}}{2908351883124559872000000} + \frac{61921208895254661989988212567311 n^{7}}{110275008901806228480000000} - \frac{9925722646170490600108095270877 n^{6}}{1934649278979056640000000} + \frac{47393917262283251844887075621 n^{5}}{1237654420895692800000} - \frac{1390548650132864765781401508097 n^{4}}{6126389383433679360000} + \frac{433846309064646404522540308733 n^{3}}{423012600284706432000} - \frac{926612557000804897407362797 n^{2}}{279770238283536000} + \frac{49062604628325366151963 n}{7220177644680} - 6644041492.
\]

\section{The count is a walk count}

The entry's defining condition is local: it is a condition on a bounded window of consecutive lines. Take as vertices the
admissible configurations of such a window and put an edge from $u$ to $v$ when $v$ may
follow $u$, which the condition decides by inspection. An admissible array is then exactly a
walk, so with $M$ the adjacency matrix of that digraph on $S=2745$ vertices, and $u,w$ the
vectors recording which windows may start and end an array,
\[
a(n)\;=\;u^{\!\top}M^{\,g(n)}w
\]
for the linear function $g$ that the entry's shape supplies. In particular $a$ is
$C$-finite: it satisfies some monic integer linear recurrence, though not necessarily a short
one.

\section{A closed form is a recurrence}

\begin{lemma}\label{lem:ann}
Let $f(m)=\sum_{j}p_j(m)\lambda_j^{\,m}$ with the $\lambda_j$ distinct and the $p_j$
polynomials, and put $q(x)=\prod_j(x-\lambda_j)^{\deg p_j+1}$. Then $q$ applied to the shift
operator annihilates $f$: writing $q(x)=x^{r}-\sum_{i=1}^{r}c_ix^{r-i}$,
\[
f(m)\;=\;\sum_{i=1}^{r}c_i\,f(m-i)\qquad\text{for every }m.
\]
\end{lemma}

\begin{proof}
The shift operator $E$ acts on $m\mapsto\lambda^{m}p(m)$ as
$(E-\lambda)\bigl(\lambda^{m}p(m)\bigr)=\lambda^{m+1}\bigl(p(m+1)-p(m)\bigr)$, and
$p(m+1)-p(m)$ has degree one less than $p$. So $(E-\lambda)^{\deg p+1}$ kills
$\lambda^{m}p(m)$, and the factors of $q$ commute.
\end{proof}

Here $q(x)=(x-1)^{32}$, of degree $32$; the closed form is a polynomial in $n$, so this is $(x-1)^{32}$, and the recurrence it encodes is
\begin{equation}\label{eq:rec}
a(n)\;=\;32\,a(n-1) - 496\,a(n-2) + 4960\,a(n-3) - 35960\,a(n-4) + 201376\,a(n-5) - 906192\,a(n-6) + 3365856\,a(n-7) - 10518300\,a(n-8) + 28048800\,a(n-9) - 64512240\,a(n-10) + 129024480\,a(n-11) - 225792840\,a(n-12) + 347373600\,a(n-13) - 471435600\,a(n-14) + 565722720\,a(n-15) - 601080390\,a(n-16) + 565722720\,a(n-17) - 471435600\,a(n-18) + 347373600\,a(n-19) - 225792840\,a(n-20) + 129024480\,a(n-21) - 64512240\,a(n-22) + 28048800\,a(n-23) - 10518300\,a(n-24) + 3365856\,a(n-25) - 906192\,a(n-26) + 201376\,a(n-27) - 35960\,a(n-28) + 4960\,a(n-29) - 496\,a(n-30) + 32\,a(n-31) - a(n-32).
\end{equation}

\section{The proof}

\begin{lemma}\label{lem:crit}
With $q$ as above, $a$ satisfies \eqref{eq:rec} for all large $n$ if and only if
$u^{\!\top}M^{\,j}q(M)w=0$ for every $j\ge0$, and it is enough to check $j<S$.
\end{lemma}

\begin{proof}
Substituting the walk expression, the residual $a(n)-\sum_ic_ia(n-i)$ equals
$u^{\!\top}M^{\,g(n)-r}q(M)w$ up to the fixed linear reparametrisation $g$. For the bound,
$M^{S}$ is an integer combination of $I,M,\dots,M^{S-1}$ by Cayley--Hamilton, so the values
$u^{\!\top}M^{j}q(M)w$ satisfy the monic recurrence given by the characteristic polynomial of
$M$; $S$ consecutive zeros therefore force all later ones, and the last nonzero value pins the
threshold exactly.
\end{proof}

\begin{theorem}
$a(n)=f(n)$ for every $n\ge11$.
\end{theorem}

\begin{proof}
The vector $w'=q(M)w$ was formed by $32$ matrix--vector products in exact integer
arithmetic, and $u^{\!\top}M^{j}w'$ evaluated for $j=0,1,\dots$, again exactly. Those integers
vanish from the point corresponding to $n=43$ onwards, and the run was continued until the
working vector was identically zero, which settles every larger $j$ at once. So $a$ satisfies
\eqref{eq:rec} for every $n>42$. By Lemma~\ref{lem:ann} $f$ satisfies \eqref{eq:rec}
identically. Both were then evaluated in exact arithmetic at the $32$ consecutive indices
$n=43,\dots,74$ and agree at each. A monic recurrence of order $32$
determines $a(m)$ from $a(m-1),\dots,a(m-32)$, and for every $m\ge75$ both
$m>42$ and $m-32\ge43$ hold, so induction on $m$ gives $a(m)=f(m)$ for every
$m\ge43$.
Below $n=43$ the recurrence argument gives nothing, since \eqref{eq:rec} is only
established for $a$ from $n=43$ on. The remaining indices $n=11,\dots,42$ were
therefore checked one at a time, directly against the walk count in exact integer arithmetic,
and agree.
\end{proof}

The range is tight in the sense that was checked: at $n=10$ the closed form and the sequence differ, so no larger range is available.

\section{Verification}

Four checks, each independent of the algebra above.

First, the walk model reproduces every one of the $21$ terms the entry publishes,
exactly, in integer arithmetic. This is what ties the model to the entry: the model is built
by reading the entry's English, and a misreading gives a different digraph and different
counts.

Second, the closed form was evaluated in exact rational arithmetic --- not floating point ---
at every published term from $n=11$ on, and agrees with the entry's own DATA at each.

Third, the annihilation test of Lemma~\ref{lem:crit} was carried out exactly, and the model
and the closed form were then compared directly at sixty consecutive indices beyond
$n=11$, far past where the proof already settles the matter.

Fourth, the closed form was deliberately perturbed --- by a constant and by a term linear in
$n$ --- and each perturbed form was rejected by the same comparison, so the test is not one
that anything would pass.

\begin{thebibliography}{9}
\bibitem{oeis} The OEIS Foundation, \emph{The On-Line Encyclopedia of Integer
Sequences}, \url{https://oeis.org}, 2026. Sequence A227331.
\bibitem{stanley} R.~P.~Stanley, \emph{Enumerative Combinatorics}, Volume 1, 2nd ed.,
Cambridge University Press, 2011. (Transfer-matrix method, Section 4.7; $C$-finite sequences,
Section 4.1.)
\bibitem{kp} M.~Kauers and P.~Paule, \emph{The Concrete Tetrahedron}, Springer, 2011.
(Closed forms and linear recurrences with constant coefficients.)
\bibitem{horn} R.~A.~Horn and C.~R.~Johnson, \emph{Matrix Analysis}, 2nd ed., Cambridge
University Press, 2013.
\end{thebibliography}

\end{document}
